% !TeX root = ../papers/nc-p3.tex
\section{指定三次式からの相手三平面の全復元}
\label{sec:marked-companion-plane-reconstruction}
\begin{theorem}\label{thm:marked-companion-plane-reconstruction}
定理\ref{thm:marked-cubic-linear-system}の固定三平面 $G'$ を取る。
$B=V(q_1,q_2)$ の非被約性も許す。$0\ne f=a_1q_1+a_2q_2$ に対し、
混合行列式が $f$ の\emph{非零倍}となる相手三平面 $F'$ は以下で全て得られる。
\begin{enumerate}
\item $g=5$ の解には平面の点 $p$ が一意に付く。$p\notin B$ では
$f(p)=0$ かつ $(a_1(p),a_2(p))\ne(0,0)$ の各点にちょうど一解があり、線型方程式で復元される。
$p\in B$ では $a_1(p)=a_2(p)=0$ の場合だけ解があり、各々非零スカラー一つで動く。
\item $g=4$ の解があるのは $f=lq$ で $q\in\langle q_1,q_2\rangle$ のときだけで、
相手はこの場合一本のアフィン直線をなす。
\item 全三次式が零となる中心はスキームとして $\{G'\}\sqcup B$ である。
特に $E=V(f)$ が滑らかなら、中心も含めた $kf$ の逆像は任意のこの $G'$ に対して
$E$ と同型な楕円正規五次曲線となる。
\end{enumerate}
\end{theorem}
\begin{proof}[核の二列と消去]
$A(v)$ を前節の三行五列の行列とし、$X(v)\in G'$ を、$G'$ での三次交代部分の
$23,31,12$ 成分を係数とするベクトルとする。三交換子が独立なので $X:V\to G'$ は同型である。
$A(v)$ は $X(v)$ と $T_v=vv^t-(v^tv)I/3$ を共に消し、$B$ の外でその二列は独立である。
従って $A$ の三次小行列式全体の双対外積は定数倍を除き $X(v)\wedge T_v$ に等しい。
実際両辺は同次数で、いずれも成分の共通零点は有限なので、有理関数の倍率は定数である。
消去二平面を $F'^\perp=\langle\alpha,\beta\rangle\subset S^*$ とすると、混合行列式は
定数倍を除いて
\[
           R(\alpha)H(\beta)-R(\beta)H(\alpha),\qquad
      R(\alpha)(v)=\alpha(X(v)),\quad H(\alpha)(v)=\alpha(T_v)
\]
となる。$R$ は $V^*$ への全射で、$H$ は $S^*$ を跡零二次式の全空間に同定する。
$\ker R$ はこの同定で $\langle q_1,q_2\rangle$ である。

$R$ が消去二平面上で階数2の場合、像を $p^\perp\subset V^*$ と書く
（ここでは点 $p$ を消す線型形式の空間）。$R$ の線型切断 $s$ と体積形式を固定する。
$p$ の局所代表から決まる二ベクトルを $l\wedge m=\iota_p\mathrm{vol}$ と分解する。
基準解 $s(l),s(m)$ の三次式を $b_1(p,v)q_1(v)+b_2(p,v)q_2(v)$ と書くと
各 $b_i$ は $p,v$ に双線型である。$v=p$ では $l,m$ が共に消えるため、
互いに素な二次式の関係から
\[
       (b_1(p,p),b_2(p,p))=\lambda(-q_2(p),q_1(p))
\]
となる。ここで $\lambda$ は定数で、零ではない。もし零なら $b_1,b_2$ は交代形式である。
切断に $q_i$ 倍の線型汎関数を加える操作は $b_i$ に $\mathrm{vol}(p,h_i,v)$ を
加えるので、両交代形式を同時に消せる。そのとき全ての $l,m$ で
$l H(s(m))=m H(s(l))$ である。三つの独立座標で割り切りを用いれば、
ある共通一次式 $L$ により $H(s(l))=lL$ となる。
全ての $l$ で $lL$ が跡零なら計量の非退化性から $L=0$ となり、切断が存在することに反する。

消去二平面を $s(l),s(m)$ のそれぞれへの鉛筆二次式の加算で全て書ける。
四つの自由係数は、各 $b_i(p,-)$ に $p$ を消す任意の一次式を加える四座標に同型に写る。
従って二つの値を $v=p$ に代入した式
\[
       z a_1(p)=-\lambda q_2(p),\qquad z a_2(p)=\lambda q_1(p)
\]
が全ての整合条件である。解 $z$ ごとに残りは一意で、$z\ne0$ の場合が第一項そのものである。
局所代表・基底を替えると $z$ は単元倍される。これは点ごとの解法であると同時に、
四自由係数の消去は各局所チャートで可逆なのでスキームの方程式でもある。
$z=0$ の全解は $B$ 自身に同型である。そこでの三平面を行列で書けば
$kv$ を保つ跡零対称行列の空間となる。$T_v\in G'$ と交換するので行列式は零であり、
$X(v)$ は $G'$ 内の交代行列の根基 $kT_v$ にあるので同じ点 $p=v$ を与える。

階数1では消去平面の基底を $q, s(l)+q'$ と取れる。三次式は $lq$ の定数倍、
$q'\in\langle q_1,q_2\rangle/kq$ が一意な自由座標である。
分解の一意性は $f=a_1q_1+a_2q_2$ の係数の一意性から従う。階数0なら $F'=G'$ だけで、
前節の同型微分により中心の単純な孤立点である。これで中心のスキーム記述も従う。

滑らかな $E$ なら $a_1,a_2$ は独立で、その共通零点 $p_*$ は $B$ の外にある。
もし $B$ 上なら $df$ もそこで消えてしまう。$p_*$ の外では二式から $z$ を正則に
消せて、ちょうど $f=0$ が残る。階数1の解はなく、$p_*$ に対応する補完点は $G'$ だけである。
$G'$ では曲線が滑らかであるから、全逆像は $E$ の滑らかな完備化として $E$ に等しい。
次数と線型正規性については、五超平面でGrassmann多様体を所期次元に切るので次数は5である。
線型包が三空間以下ならその中で中心超平面が切る長さ5の零スキームが平面以下に入る。
より直接に、Grassmann多様体のPl\"ucker環はCohen--Macaulayで、
所期余次元の五つの一次式は正則列となる。商は深さ2で飽和し、残る一次部分の次元は5である。
従って次数五の完全線型系による埋込みである。
\path{computations/verify_companion_plane_reconstruction.py} は四基点が別々の場合と
単一局所環 $k[y]/(y^4)$ の場合に、核の二列と $\lambda\ne0$ の恒等式を検算する。
\end{proof}
